\documentclass[10pt,a4paper]{amsart}
\usepackage[margin=19mm]{geometry}
\usepackage{amsmath,amssymb,mathtools}
\usepackage[scr=rsfs,cal=euler]{mathalfa}
\usepackage{xfrac}
\usepackage[unicode,hidelinks,bookmarks=false]{hyperref}
\newcommand{\ie}{i.e.\ }
\newcommand{\cf}{cf.\ }
\newcommand{\resp}{respectively\ }
\newcommand{\Subsection}{\subsection}
\providecommand{\nobreakdash}{\nobreak\hbox{-}}
\setlength{\emergencystretch}{2em}
\multlinegap0pt
\usepackage{bm}
\usepackage{bbm}
\usepackage{dsfont}
\usepackage{xspace}
\usepackage{cancel}
\usepackage{mathtools}
\usepackage{amsthm}
\makeatletter
\@ifundefined{theorem}{\newtheorem{theorem}{Theorem}}{}
\@ifundefined{lemma}{\newtheorem{lemma}[theorem]{Lemma}}{}
\makeatother
\newcommand{\N}{\mathbb{N}}
\newcommand{\lpf}{\normalfont\text{lpf}}
\newcommand{\mO}{\mathcal{O}}
\newcommand{\p}{\mathfrak{p}}
\renewcommand{\tilde}{\widetilde}
\title[Polynomial actions of rings of integers of global fields and]{Polynomial actions of rings of integers of global fields and quasirandomness of Paley-type graphs}
\author{Ethan Ackelsberg, Vitaly Bergelson}
\date{Source: arXiv:2509.17868v1 (CC BY 4.0)}
\begin{document}
\maketitle

\noindent\emph{Source and licence.} Polynomial actions of rings of integers of global fields and quasirandomness of Paley-type graphs. Version arXiv:2509.17868v1 (CC BY 4.0), distributed under the Creative Commons Attribution 4.0 International licence, as recorded in the arXiv licence metadata for this exact version and shown on the abstract page https://arxiv.org/abs/2509.17868v1. Full rights notice: licenses/arxiv-2509.17868-CC-BY-4.0.txt.\par\smallskip

\noindent\textit{Editorial scope.} Counted unit: the Lemma whose source label is \texttt{lem: nonzero polynomial} in the article (source0.tex lines 2042--2048, source label \texttt{lem: nonzero polynomial}), delivered together with the article's own complete original proof of it (source0.tex lines 2050--2109, one \texttt{proof} environment of 60 lines). No mathematics is written, repaired or re-derived: statement and proof are the article's own text line for line. Required context retained uncounted: none. Every definition, display and label of those lines is retained unchanged. Omitted from this package and not redistributed: 1--2041, 2049--2049, 2110--2776. Nothing inside a retained excerpt is omitted or altered: every retained body is delivered exactly as the article's lines are written, so no omission receipt of any kind accompanies this package. Citations: the retained excerpts carry no citation command at all, so no bibliography entry of the article is transcribed and no reference is redistributed. Reference adaptation: none was needed and none was made; every reference inside the delivered unit resolves inside it, so no unresolved reference or citation is produced. Printed slips kept literal: no printed word, symbol, hypothesis or number of the counted statement or of its proof was altered. Layout and class: the article's own class is \texttt{amsart} with packages including geometry, graphicx, amsmath, amssymb, amsfonts, mathrsfs, amsthm, hyperref, enumerate, bbm, bm, tikz, makecell, floatrow; the wrapper is the lane's pinned \texttt{amsart} editorial wrapper at 10pt on A4, which restates the theorem-like environments the retained text needs and adds the article's own macro definitions that the retained text needs (\texttt{\textbackslash N}, \texttt{\textbackslash lpf}, \texttt{\textbackslash mO}, \texttt{\textbackslash p}, \texttt{\textbackslash tilde}), with extra wrapper packages (bm, bbm, dsfont, xspace, cancel, mathtools). The delivered numbering is the unit's own. Assets: no retained span contains a figure environment, a graphics-inclusion command, a PDF-page inclusion, a table or a TikZ picture; no image file is copied into this package, the package declares no asset dependency and no image file is redistributed with it. Licence: version arXiv:2509.17868v1 of this article is distributed under the Creative Commons Attribution 4.0 International licence, as recorded in the arXiv licence metadata for this exact version and shown on the abstract page https://arxiv.org/abs/2509.17868v1. A full rights notice is delivered at licenses/arxiv-2509.17868-CC-BY-4.0.txt. Characters: every retained excerpt is pure ASCII, so the delivered engine, XeLaTeX with its default Latin Modern text font, reports no missing character.
	\begin{lemma} \label{lem: nonzero polynomial}
		Let $l \in \N$, and let $T(x_1, \dots, x_l)$ be a nonzero polynomial with coefficients in $I^{\perp}$ of degree $d_i$ in the variable $x_i$ for $i = 1, \dots, l$.
		Then
		\begin{equation*}
			\left| \left\{ (x_1, \dots, x_l) \in (\mO_K/I)^l : T(x_1, \dots, x_l) = 0 \pmod{\mO_K} \right\} \right| \le \left( \sum_{i=1}^l{d_i} \right) \frac{[\mO_K:I]^l}{\lpf(I)}.
		\end{equation*}
	\end{lemma}

	\begin{proof}[Proof of Lemma]
		Let us first consider the case $l = 1$.
        Write $T(x) = \alpha_d x^d + \ldots + \alpha_1 x + \alpha_0$ with $\alpha_0, \ldots, \alpha_d \in I^{\perp}$.
        We view $T$ as a map on $\mO_K$ with $T(I) \subseteq \mO_K$.
        Then let $J$ be the ideal
        \begin{equation*}
            \left( \textup{span}_{\mO_K}\left\{ \alpha_0, \ldots, \alpha_d \right\} \right)^{\perp}.
        \end{equation*}
        In other words, $J$ is the largest ideal such that $T(J) \subseteq \mO_K$.
        By construction, $J \supseteq I$ and $J$ is a proper ideal of $\mO_K$.
        Let $\p$ be a prime ideal with $\p \mid J$.
        Letting $\tilde{T}$ be the reduction of $T$ mod $\p$, we have that $\tilde{T}$ is a nonzero polynomial.
        But $\mO_K/\p$ is a field, so $\tilde{T}$ can have at most $d$ roots mod $\p$.
		
		Now suppose $T(x) = 0 \pmod{\mO_K}$.
		Then $\tilde{T}(x) = 0$.
		Therefore,
		\begin{equation*}
			\left| \{x \in \mO_K/I : T(x) = 0 \pmod{\mO_K} \} \right| \le d \frac{[\mO_K:I]}{[\mO_K:\p]} \le d \frac{[\mO_K:I]}{\lpf(I)}.
		\end{equation*} \\
		
		Suppose $l \ge 2$.
		If $\lpf(I) \le \sum_{i=1}^l{d_i}$, then there is nothing to prove, so assume $\lpf(I) > \sum_{i=1}^l{d_i}$.
		Fix $x \in \mO_K/I$, and let $T_x(x_1, \dots, x_{l-1}) = T(x_1, \dots, x_{l-1}, x)$.
		If $T_x(x_1, \dots, x_{l-1})$ is not the zero polynomial, then by the induction hypothesis,
		\begin{equation*}
			\left| \left\{ (x_1, \dots, x_{l-1}) \in (\mO_K/I)^{l-1} : T_x(x_1, \dots, x_{l-1}) = 0 \right\} \right| \le \left( \sum_{i=1}^{l-1}{d_i} \right) \frac{[\mO_K:I]^{l-1}}{\lpf(I)}.
		\end{equation*}
		Hence,
		\begin{align*}
			& \left| \left\{ (x_1, \dots, x_l) \in (\mO_K/I)^l : T(x_1, \dots, x_l) = 0 \right\} \right| \\
			 &~= [\mO_K:I]^{l-1} \left| \{x \in \mO_K/I : T_x = 0\} \right|
			 + \left| \{(x_1, \dots, x_{l-1}, x) : T_x \ne 0, T_x(x_1, \dots, x_{l-1}) = 0\} \right| \\
			 &~\le [\mO_K:I]^{l-1} \left| \{x \in \mO_K/I : T_x = 0\} \right| + \left( \sum_{i=1}^{l-1}{d_i} \right) \frac{[\mO_K:I]^l}{\lpf(I)}.
		\end{align*}
		It therefore suffices to prove
		\begin{equation*}
			\left| \{x \in \mO_K/I : T_x = 0\} \right| \le d_l \frac{[\mO_K:I]}{\lpf(I)}.
		\end{equation*}
		
		Fix $x_1, \dots, x_{l-1} \in \mO_K/I$, and let $T^{x_1, \dots, x_{l-1}}(x) = T_x(x_1, \dots, x_{l-1}) = T(x_1, \dots, x_{l-1}, x)$.
		If $T_x = 0$, then $T^{x_1, \dots, x_{l-1}}(x) = 0$.
		Therefore, by the $l=1$ case above, it follows that
		\begin{equation*}
			\left| \{x \in \mO_K/I : T_x = 0\} \right|
			 \le \left| \{x \in \mO_K/I : T^{x_1, \dots, x_{l-1}}(x) = 0\} \right|
			 \le d_l \frac{[\mO_K:I]}{\lpf(I)},
		\end{equation*}
		unless $T^{x_1, \dots, x_{l-1}}$ is the zero polynomial.
		So, it remains to find $x_1, \dots, x_{l-1}$ such that $T^{x_1, \dots, x_{l-1}}$ is a nonzero polynomial.
		Note that the coefficients of $T^{x_1, \dots, x_{l-1}}$ are polynomial expressions in $x_1, \dots, x_{l-1}$ of degree at most $d_i$ in the variable $x_i$.
		Since $T$ is not the zero polynomial, there is at least one coefficient that is a nonzero polynomial $C(x_1, \dots, x_{l-1})$.
		By the induction hypothesis,
		\begin{equation*}
			\left| \left\{ (x_1, \dots, x_{l-1}) \in (\mO_K/I)^{l-1} : C(x_1, \dots, x_{l-1}) = 0 \right\} \right|
			 \le \left( \sum_{i=1}^{l-1}{d_i} \right) \frac{[\mO_K:I]^{l-1}}{\lpf(I)}.
		\end{equation*}
		Since $\lpf(I) > \sum_{i=1}^l{d_i} \ge \sum_{i=1}^{l-1}{d_i}$ by assumption, it follows that $C(x_1, \dots, x_{l-1}) \ne 0$ for some $(x_1, \dots, x_{l-1}) \in (\mO_K/I)^{l-1}$.
		For this choice of $x_1, \dots, x_{l-1}$, the polynomial $T^{x_1, \dots, x_{l-1}}(x)$ is not the zero polynomial, so we are done.
	\end{proof}

\end{document}
